% Generated draft: requires mathematical and source review.
\documentclass[11pt]{article}
\usepackage[margin=25mm]{geometry}
\usepackage{fontspec}
\setmainfont{Latin Modern Roman}
\usepackage{amsmath,amssymb,mathtools,unicode-math}
\setmathfont{Latin Modern Math}
\usepackage{enumitem,needspace,xurl,hyperref}
\setlength{\parindent}{0pt}
\setlength{\parskip}{5pt}
\setlength{\emergencystretch}{3em}
\newcommand{\N}{\mathbb{N}}
\newcommand{\Z}{\mathbb{Z}}
\newcommand{\Q}{\mathbb{Q}}
\newcommand{\R}{\mathbb{R}}
\newcommand{\C}{\mathbb{C}}
\newcommand{\F}{\mathbb{F}}
\newcommand{\A}{\mathbb{A}}
\newcommand{\PP}{\mathbb{P}}
\newcommand{\E}{\mathbb{E}}
\newcommand{\eps}{\varepsilon}
\newcommand{\dd}{\,\mathrm{d}}
\begin{document}
\section*{UnsolvedMath ID: 800244}
\textbf{Source problem code:} catalogue entry
\subsection{Definitions and mathematical statement}
The supplied question concerns a smooth pseudoconvex domain and compares its regular order of contact with its commutator type, with the qualification ``one vector field at a time.'' The record does not define these invariants or explain that qualification.

To describe the usual local terminology, let \(\Omega\subset\C^n\) be a domain with \(C^\infty\) boundary, let \(p\in\partial\Omega\), and choose a real smooth defining function \(\rho\) near \(p\), so that \(\Omega=\{\rho<0\}\) locally and \(\dd\rho\ne0\) on the boundary. The complex tangent space is
\[
T_q^{1,0}\partial\Omega=\left\{v\in\C^n:\sum_{j=1}^n\rho_{z_j}(q)v_j=0\right\}.
\]
Pseudoconvexity means that the Levi form \(\sum_{j,k}\rho_{z_j\bar z_k}(q)v_j\overline{v_k}\) is nonnegative on this space at each boundary point \(q\).

One standard scalar meaning of regular order of contact at \(p\) is
\[
\Delta_{\mathrm{reg}}(\partial\Omega,p)
=\sup_{\substack{\gamma:(\C,0)\to(\C^n,p)\ \mathrm{holomorphic\ germ}\\\gamma'(0)\ne0}}
\operatorname{ord}_0(\rho\circ\gamma).
\]
Here the order of a smooth germ \(f(t,\bar t)\) vanishing at zero is the least \(a+b\) for which \(\partial_t^a\partial_{\bar t}^b f(0)\ne0\), and is \(+\infty\) if all such derivatives vanish.

A usual single-field commutator construction starts with a smooth local section \(L\) of \(T^{1,0}\partial\Omega\) with \(L(p)\ne0\). Put \(X=L+\bar L\), \(Y=i(L-\bar L)\), and let
\[
\theta=\left.\frac{i}{2}(\bar\partial\rho-\partial\rho)\right|_{T\partial\Omega}.
\]
The associated bracket length is
\[
c(L,p)=\inf\left\{m\ge2:\theta_p(B(p))\ne0\text{ for some iterated Lie bracket }B\text{ of length }m\text{ in }X,Y\right\},
\]
with \(\inf\varnothing=+\infty\). Length counts the occurrences of \(X\) and \(Y\).

\textit{Source question, attributed to J. McNeal.} Is the regular order of contact of a smooth pseudoconvex domain the same as the commutator type, ``one vector field at a time''?

\textit{Unresolved formulation.} The definitions above explain standard terminology; the record does not establish that they are the intended conventions. In particular, the scalar quantity \(\Delta_{\mathrm{reg}}(\partial\Omega,p)\) does not itself select a vector field. A precise equality requires a definition of regular contact associated with an individual field or direction, the intended commutator convention, and the quantifiers over boundary points and fields. These choices cannot be recovered uniquely from the supplied text. The ambient dimension, any finite-type hypothesis, and the treatment of infinite orders are also unspecified.
\subsection{Short English statement}
J. McNeal asks whether regular holomorphic contact and commutator type agree for a smooth pseudoconvex domain, considered one vector field at a time. The supplied record does not specify the directional contact invariant or the exact quantifiers needed for a definite equality. Its \texttt{open} label is supplied registry metadata, not an independently verified status.
\subsection{Sources}
\begin{itemize}
\item Local UnsolvedMath registry, permanent ID \(800244\), exact source code \texttt{catalogue entry}, title \emph{Several Complex Variables Problem List --- Problem 42}. The supplied record labels the problem \texttt{open}.
\item Buzzard, Fornaess, and Krantz, \emph{The Several Complex Variables Problem List} (1995), Problem 42. Supplied source URL: \url{https://arxiv.org/abs/math/9509202}. The record reports access on 2026-10-16 and extraction by \texttt{arxiv\_latex}; these are supplied provenance details, not a claim of independent source access or verification.
\item The question itself is attributed in the supplied statement to J. McNeal, with contact address \texttt{mcneal@math.princeton.edu}. The registry's proposer field names Buzzard, Fornaess, and Krantz and gives the year 1995; this collection-level attribution is distinct from the question's attribution.
\end{itemize}
\end{document}
